\documentclass[11pt,a4paper]{article}

% ------------------------------------------------------------
% Packages
% ------------------------------------------------------------
\usepackage[margin=1in]{geometry}
\usepackage{amsmath,amssymb,amsthm,mathtools}
\usepackage{booktabs,longtable,array,tabularx}
\usepackage{enumitem}
\usepackage{xcolor}
\usepackage{hyperref}
\usepackage{titlesec}
\usepackage{fancyhdr}
\usepackage{multicol}

% ------------------------------------------------------------
% Formatting
% ------------------------------------------------------------
\hypersetup{
    colorlinks=true,
    linkcolor=blue!60!black,
    urlcolor=blue!60!black,
    citecolor=blue!60!black
}

\setlist[itemize]{leftmargin=*,noitemsep,topsep=2pt}
\setlist[enumerate]{leftmargin=*,noitemsep,topsep=2pt}

\pagestyle{fancy}
\fancyhf{}
\rhead{Time Series Analysis}
\lhead{Course Structure}
\cfoot{\thepage}

\titleformat{\section}
  {\Large\bfseries}{\thesection}{0.7em}{}
\titleformat{\subsection}
  {\large\bfseries}{\thesubsection}{0.7em}{}

% ------------------------------------------------------------
% Custom commands
% ------------------------------------------------------------
\newcommand{\AR}{\operatorname{AR}}
\newcommand{\MA}{\operatorname{MA}}
\newcommand{\ARMA}{\operatorname{ARMA}}
\newcommand{\ARIMA}{\operatorname{ARIMA}}
\newcommand{\GARCH}{\operatorname{GARCH}}
\newcommand{\ARCH}{\operatorname{ARCH}}
\newcommand{\TAR}{\operatorname{TAR}}
\newcommand{\E}{\mathbb{E}}
\newcommand{\Var}{\operatorname{Var}}
\newcommand{\Cov}{\operatorname{Cov}}
\newcommand{\Corr}{\operatorname{Corr}}

% ============================================================
\begin{document}

\begin{center}
    {\Huge\bfseries Time Series Analysis}\\[0.4cm]
    {\Large Full Course Structure, Modules, Concepts, Learning Objectives}\\[0.4cm]
    {\large Based on the supplied textbook Table of Contents}
\end{center}

\vspace{0.5cm}

% ============================================================
\section{Course Overview}

\textbf{Course Title:} Time Series Analysis

\textbf{Course Level:} Advanced undergraduate / Graduate

\textbf{Primary Mathematical Orientation:}
Probability, mathematical statistics, regression, stochastic processes, and
statistical modeling.

\textbf{Primary Computational Orientation:}
Statistical computing and empirical time-series analysis, with particular
emphasis on implementation in R.

\subsection{Course Description}

This course develops a systematic framework for analyzing observations
collected sequentially through time. The course begins with the fundamental
concepts of stochastic processes, means, variances, covariances, and
stationarity. It then develops deterministic and stochastic trend models,
stationary AR, MA, and ARMA models, and nonstationary ARIMA models.

The central part of the course emphasizes the complete time-series
model-building cycle:

\[
\boxed{
\text{Identification}
\rightarrow
\text{Specification}
\rightarrow
\text{Estimation}
\rightarrow
\text{Diagnostics}
\rightarrow
\text{Forecasting}
}
\]

The course subsequently extends these ideas to seasonal models, time-series
regression, heteroscedastic models such as ARCH and GARCH, spectral analysis,
and nonlinear threshold models.

\subsection{Overall Course Goals}

By the end of the course, students should be able to:

\begin{enumerate}
    \item Formulate real-world sequential data as stochastic time-series
          models.
    \item Distinguish deterministic trends from stochastic trends.
    \item Understand and assess stationarity.
    \item Derive and interpret autocovariance and autocorrelation structures.
    \item Analyze MA, AR, and ARMA processes.
    \item Model nonstationary data using differencing and ARIMA models.
    \item Identify plausible models using ACF, PACF, and related diagnostics.
    \item Estimate model parameters using several statistical methods.
    \item Diagnose model adequacy using residual analysis.
    \item Produce and evaluate multi-step forecasts.
    \item Model seasonal time-series behavior.
    \item Incorporate interventions, outliers, and explanatory variables.
    \item Model conditional heteroscedasticity using ARCH/GARCH methods.
    \item Analyze time series in the frequency domain.
    \item Recognize and model nonlinear dynamics using threshold models.
    \item Implement time-series methods computationally, particularly using R.
\end{enumerate}

% ============================================================
\section{Prerequisite Knowledge}

The following prerequisites are recommended before beginning the course.

\subsection{Probability}

Students should be comfortable with:

\begin{itemize}
    \item Random variables.
    \item Probability distributions.
    \item Expectation.
    \item Variance.
    \item Covariance.
    \item Correlation.
    \item Conditional expectation.
    \item Basic properties of independent and identically distributed random
          variables.
\end{itemize}

\subsection{Mathematical Statistics}

Students should know:

\begin{itemize}
    \item Point estimation.
    \item Sampling distributions.
    \item Least squares.
    \item Maximum likelihood.
    \item Statistical hypothesis testing.
    \item Confidence intervals.
    \item Bias and variance.
    \item Consistency and efficiency at an introductory level.
\end{itemize}

\subsection{Regression}

Students should understand:

\begin{itemize}
    \item Simple and multiple linear regression.
    \item Regression coefficients.
    \item Residuals.
    \item Standard errors.
    \item $t$-tests and $F$-tests.
    \item Model fit.
    \item Residual diagnostics.
\end{itemize}

\subsection{Linear Algebra and Calculus}

A working knowledge of:

\begin{itemize}
    \item Vectors and matrices.
    \item Matrix multiplication.
    \item Eigenvalues and eigenvectors at a basic level.
    \item Differentiation.
    \item Optimization.
    \item Basic infinite series.
\end{itemize}

\subsection{Computational Prerequisites}

Students should be able to:

\begin{itemize}
    \item Import and manipulate data.
    \item Produce basic statistical plots.
    \item Fit statistical models.
    \item Interpret numerical output.
    \item Write basic R code.
\end{itemize}

% ============================================================
\section{Course Architecture}

The 15 textbook chapters can be organized into the following pedagogical
modules.

\begin{longtable}{p{0.10\textwidth}p{0.25\textwidth}p{0.55\textwidth}}
\toprule
\textbf{Module} & \textbf{Theme} & \textbf{Main Chapters} \\
\midrule
1 & Foundations &
Chapters 1--2: Introduction and Fundamental Concepts \\

2 & Trends and Regression &
Chapter 3: Trends \\

3 & Stationary Models &
Chapter 4: Models for Stationary Time Series \\

4 & Nonstationary Models &
Chapter 5: Models for Nonstationary Time Series \\

5 & Model Identification &
Chapter 6: Model Specification \\

6 & Parameter Estimation &
Chapter 7: Parameter Estimation \\

7 & Model Diagnostics &
Chapter 8: Model Diagnostics \\

8 & Forecasting &
Chapter 9: Forecasting \\

9 & Seasonal Models &
Chapter 10: Seasonal Models \\

10 & Time-Series Regression &
Chapter 11: Time Series Regression Models \\

11 & Volatility Modeling &
Chapter 12: Time Series Models of Heteroscedasticity \\

12 & Frequency-Domain Methods &
Chapters 13--14: Spectral Analysis and Spectrum Estimation \\

13 & Nonlinear Models &
Chapter 15: Threshold Models \\

14 & Statistical Computing &
R commands, TSA library, datasets, and computational implementation \\
\bottomrule
\end{longtable}

% ============================================================
\section{Module 1: Introduction and Fundamental Concepts}

\subsection{Chapter 1 -- Introduction}

\textbf{Major Topics}

\begin{itemize}
    \item Examples of time series.
    \item Characteristics of temporally ordered observations.
    \item Model-building strategy.
    \item Historical development of time-series plots.
    \item Overview of time-series methodology.
\end{itemize}

\textbf{Core Concepts}

\begin{itemize}
    \item Time-indexed observations.
    \item Dependence across time.
    \item Trend.
    \item Random variation.
    \item Model identification.
    \item Model estimation.
    \item Model checking.
    \item Forecasting.
\end{itemize}

\textbf{Learning Objectives}

Students will be able to:

\begin{enumerate}
    \item Define a time series and distinguish it from ordinary cross-sectional
          data.
    \item Identify common features of time-series data.
    \item Explain why temporal dependence changes statistical analysis.
    \item Describe the basic model-building strategy.
    \item Interpret time-series plots as a first diagnostic tool.
\end{enumerate}

\subsection{Chapter 2 -- Fundamental Concepts}

\textbf{Major Topics}

\begin{itemize}
    \item Time series and stochastic processes.
    \item Means.
    \item Variances.
    \item Covariances.
    \item Stationarity.
\end{itemize}

\textbf{Core Mathematical Concepts}

For a stochastic process $\{X_t\}$:

\[
\mu_t = \E(X_t),
\]

\[
\gamma(t,s)
=
\Cov(X_t,X_s),
\]

and the autocorrelation function may be expressed as

\[
\rho(t,s)
=
\frac{\Cov(X_t,X_s)}
{\sqrt{\Var(X_t)\Var(X_s)}}.
\]

For a weakly stationary process,

\[
\E(X_t)=\mu,
\]

\[
\Var(X_t)=\gamma(0),
\]

and

\[
\Cov(X_t,X_{t-k})=\gamma(k),
\]

where the covariance depends on the lag $k$ rather than on calendar time
itself.

\textbf{Learning Objectives}

Students will be able to:

\begin{enumerate}
    \item Define a stochastic process.
    \item Distinguish a stochastic process from a single observed time series.
    \item Calculate means, variances, and covariances.
    \item Explain autocovariance and autocorrelation.
    \item Distinguish stationary and nonstationary processes.
    \item Explain why stationarity is central to classical time-series modeling.
\end{enumerate}

\textbf{Difficult Topics}

\begin{itemize}
    \item The distinction between a stochastic process and its realization.
    \item Stationarity.
    \item Lag-based covariance structures.
    \item Interpretation of dependence over time.
\end{itemize}

% ============================================================
\section{Module 2: Trends and Regression}

\subsection{Chapter 3 -- Trends}

\textbf{Major Topics}

\begin{itemize}
    \item Deterministic versus stochastic trends.
    \item Estimation of a constant mean.
    \item Regression methods.
    \item Reliability and efficiency of regression estimates.
    \item Interpretation of regression output.
    \item Residual analysis.
\end{itemize}

\textbf{Core Concepts}

A deterministic trend may be represented by

\[
X_t = \beta_0+\beta_1t+\epsilon_t,
\]

whereas a stochastic trend may arise through accumulated random shocks,
for example,

\[
X_t=X_{t-1}+\epsilon_t.
\]

\textbf{Learning Objectives}

Students will be able to:

\begin{enumerate}
    \item Distinguish deterministic and stochastic trends.
    \item Estimate trend parameters using regression.
    \item Interpret regression coefficients for time-indexed data.
    \item Evaluate the reliability of regression estimates.
    \item Analyze residuals for evidence of model inadequacy.
    \item Explain why ordinary regression assumptions may fail for time-series
          observations.
\end{enumerate}

\textbf{Difficult Topics}

\begin{itemize}
    \item Deterministic versus stochastic trends.
    \item Efficiency of regression estimators under temporal dependence.
    \item Residual analysis for dependent observations.
\end{itemize}

% ============================================================
\section{Module 3: Stationary Time-Series Models}

\subsection{Chapter 4 -- Models for Stationary Time Series}

\textbf{Major Topics}

\begin{itemize}
    \item General linear processes.
    \item Moving average processes.
    \item Autoregressive processes.
    \item Mixed autoregressive moving average models.
    \item Invertibility.
\end{itemize}

\subsection{Moving Average Models}

An $\MA(q)$ process can be written as

\[
X_t
=
\mu+
\epsilon_t
+
\theta_1\epsilon_{t-1}
+\cdots+
\theta_q\epsilon_{t-q}.
\]

Students study:

\begin{itemize}
    \item Finite dependence.
    \item Mean and variance.
    \item Autocovariance.
    \item ACF behavior.
    \item Parameter interpretation.
\end{itemize}

\subsection{Autoregressive Models}

An $\AR(p)$ process can be represented by

\[
X_t
=
c+
\phi_1X_{t-1}
+\cdots+
\phi_pX_{t-p}
+\epsilon_t.
\]

Important concepts include:

\begin{itemize}
    \item Stability.
    \item Stationarity regions.
    \item Characteristic equations.
    \item Autocorrelation structure.
\end{itemize}

\subsection{ARMA Models}

The mixed model combines autoregressive and moving-average components:

\[
\phi(B)X_t
=
\theta(B)\epsilon_t,
\]

where $B$ denotes the backshift operator.

\textbf{Learning Objectives}

Students will be able to:

\begin{enumerate}
    \item Explain the structure of linear stochastic processes.
    \item Derive basic properties of MA processes.
    \item Analyze stationarity conditions for AR processes.
    \item Construct and interpret ARMA models.
    \item Understand the distinction between stationarity and invertibility.
    \item Relate model structure to autocorrelation behavior.
\end{enumerate}

\textbf{Difficult Topics}

\begin{itemize}
    \item Stationarity regions for AR models.
    \item Characteristic roots.
    \item ARMA representation.
    \item Invertibility.
    \item Derivation of autocorrelation functions.
\end{itemize}

% ============================================================
\section{Module 4: Nonstationary Time-Series Models}

\subsection{Chapter 5 -- Models for Nonstationary Time Series}

\textbf{Major Topics}

\begin{itemize}
    \item Stationarity through differencing.
    \item ARIMA models.
    \item Constant terms in ARIMA models.
    \item Transformations.
\end{itemize}

\subsection{Differencing}

The first difference is

\[
\nabla X_t=X_t-X_{t-1}.
\]

Higher-order differencing is represented as

\[
\nabla^d X_t=(1-B)^dX_t.
\]

\subsection{ARIMA Models}

An $\ARIMA(p,d,q)$ model can be represented as

\[
\phi(B)(1-B)^dX_t
=
\theta(B)\epsilon_t.
\]

\textbf{Learning Objectives}

Students will be able to:

\begin{enumerate}
    \item Recognize common forms of nonstationarity.
    \item Explain why differencing can produce stationarity.
    \item Select an appropriate order of differencing conceptually.
    \item Interpret the components $p$, $d$, and $q$ of an ARIMA model.
    \item Understand the role of constant terms in ARIMA models.
    \item Apply appropriate transformations when required.
\end{enumerate}

\textbf{Difficult Topics}

\begin{itemize}
    \item Distinguishing trend from stochastic nonstationarity.
    \item Selecting the order of differencing.
    \item Understanding ARIMA notation.
    \item Backshift-operator algebra.
\end{itemize}

% ============================================================
\section{Module 5: Model Specification and Identification}

\subsection{Chapter 6 -- Model Specification}

\textbf{Major Topics}

\begin{itemize}
    \item Properties of the sample autocorrelation function.
    \item Partial autocorrelation function.
    \item Extended autocorrelation function.
    \item Specification using simulated series.
    \item Specification under nonstationarity.
    \item Alternative specification methods.
    \item Specification using actual time series.
\end{itemize}

\subsection{Autocorrelation Function}

The sample ACF at lag $k$ is conceptually based on

\[
r_k
=
\frac{
\sum_{t=k+1}^{n}(X_t-\bar X)(X_{t-k}-\bar X)
}{
\sum_{t=1}^{n}(X_t-\bar X)^2
}.
\]

\subsection{Partial Autocorrelation}

The PACF measures the relationship between $X_t$ and $X_{t-k}$ after
accounting for intermediate lags.

\textbf{Learning Objectives}

Students will be able to:

\begin{enumerate}
    \item Compute and interpret the sample ACF.
    \item Explain the purpose of the PACF.
    \item Use ACF/PACF patterns to propose ARMA models.
    \item Recognize characteristic patterns associated with AR and MA models.
    \item Account for nonstationarity during model identification.
    \item Compare alternative candidate specifications.
\end{enumerate}

\textbf{Difficult Topics}

\begin{itemize}
    \item Correct interpretation of ACF and PACF patterns.
    \item Distinguishing theoretical and sample correlation behavior.
    \item Model identification under nonstationarity.
    \item Avoiding over-identification and under-identification.
\end{itemize}

% ============================================================
\section{Module 6: Parameter Estimation}

\subsection{Chapter 7 -- Parameter Estimation}

\textbf{Major Topics}

\begin{itemize}
    \item Method of moments.
    \item Least-squares estimation.
    \item Maximum likelihood.
    \item Unconditional least squares.
    \item Properties of parameter estimates.
    \item Parameter-estimation illustrations.
    \item Bootstrapping ARIMA models.
\end{itemize}

\textbf{Learning Objectives}

Students will be able to:

\begin{enumerate}
    \item Explain the method of moments.
    \item Estimate ARIMA parameters using least-squares approaches.
    \item Explain the basic principle of maximum likelihood estimation.
    \item Compare estimation procedures.
    \item Interpret estimated parameters and their uncertainty.
    \item Understand the role of bootstrap procedures for ARIMA models.
\end{enumerate}

\textbf{Difficult Topics}

\begin{itemize}
    \item Likelihood construction for dependent observations.
    \item Conditional versus unconditional estimation.
    \item Statistical properties of estimators.
    \item Bootstrap procedures for dependent data.
\end{itemize}

% ============================================================
\section{Module 7: Model Diagnostics}

\subsection{Chapter 8 -- Model Diagnostics}

\textbf{Major Topics}

\begin{itemize}
    \item Residual analysis.
    \item Overfitting.
    \item Parameter redundancy.
\end{itemize}

A fitted model should produce residuals that behave approximately like
uncorrelated random errors:

\[
\hat{\epsilon}_t
=
X_t-\hat X_t.
\]

\textbf{Learning Objectives}

Students will be able to:

\begin{enumerate}
    \item Compute and inspect model residuals.
    \item Determine whether residual autocorrelation remains.
    \item Assess whether a fitted model adequately captures temporal dependence.
    \item Identify overfitting.
    \item Recognize redundant parameters.
    \item Decide whether a model should be revised.
\end{enumerate}

\textbf{Difficult Topics}

\begin{itemize}
    \item Distinguishing random residual behavior from remaining structure.
    \item Understanding overfitting.
    \item Balancing model complexity and adequacy.
\end{itemize}

% ============================================================
\section{Module 8: Forecasting}

\subsection{Chapter 9 -- Forecasting}

\textbf{Major Topics}

\begin{itemize}
    \item Minimum mean-square-error forecasting.
    \item Forecasting deterministic trends.
    \item ARIMA forecasting.
    \item Prediction limits.
    \item Forecasting illustrations.
    \item Updating forecasts.
    \item Forecast weights.
    \item Exponentially weighted moving averages.
    \item Forecasting transformed series.
\end{itemize}

\subsection{Minimum Mean-Square-Error Forecast}

For a future value $X_{t+h}$, the MMSE forecast is

\[
\hat X_t(h)
=
\E(X_{t+h}\mid \mathcal{F}_t),
\]

where $\mathcal{F}_t$ denotes the information available at time $t$.

\textbf{Learning Objectives}

Students will be able to:

\begin{enumerate}
    \item Explain the principle of minimum mean-square-error forecasting.
    \item Construct one-step and multi-step forecasts.
    \item Forecast using ARIMA models.
    \item Construct prediction limits.
    \item Update forecasts as new observations become available.
    \item Explain forecast weights.
    \item Apply exponentially weighted moving-average ideas.
    \item Account for transformations when forecasting.
\end{enumerate}

\textbf{Difficult Topics}

\begin{itemize}
    \item Conditional expectation.
    \item Multi-step forecasting.
    \item Forecast-error variance.
    \item Prediction intervals.
    \item Forecasting transformed series.
\end{itemize}

% ============================================================
\section{Module 9: Seasonal Time-Series Models}

\subsection{Chapter 10 -- Seasonal Models}

\textbf{Major Topics}

\begin{itemize}
    \item Seasonal ARIMA models.
    \item Multiplicative seasonal ARMA models.
    \item Nonstationary seasonal ARIMA models.
    \item Seasonal model specification.
    \item Model fitting and checking.
    \item Seasonal forecasting.
\end{itemize}

For seasonal period $s$, seasonal differencing is represented by

\[
\nabla_s X_t
=
X_t-X_{t-s}.
\]

A general seasonal ARIMA structure combines ordinary and seasonal operators.

\textbf{Learning Objectives}

Students will be able to:

\begin{enumerate}
    \item Identify seasonal patterns in time-series data.
    \item Explain seasonal differencing.
    \item Construct seasonal ARIMA specifications.
    \item Understand multiplicative seasonal structures.
    \item Fit and diagnose seasonal models.
    \item Produce forecasts from seasonal models.
\end{enumerate}

\textbf{Difficult Topics}

\begin{itemize}
    \item Seasonal versus nonseasonal dependence.
    \item Seasonal differencing.
    \item Multiplicative seasonal ARIMA notation.
    \item Seasonal model identification.
\end{itemize}

% ============================================================
\section{Module 10: Time-Series Regression Models}

\subsection{Chapter 11 -- Time Series Regression Models}

\textbf{Major Topics}

\begin{itemize}
    \item Intervention analysis.
    \item Outliers.
    \item Spurious correlation.
    \item Prewhitening.
    \item Stochastic regression.
\end{itemize}

\textbf{Learning Objectives}

Students will be able to:

\begin{enumerate}
    \item Incorporate interventions into time-series models.
    \item Identify and account for outliers.
    \item Explain the danger of spurious correlation.
    \item Understand prewhitening.
    \item Distinguish ordinary regression from stochastic regression.
    \item Model relationships between multiple time-dependent variables.
\end{enumerate}

\textbf{Difficult Topics}

\begin{itemize}
    \item Spurious correlation.
    \item Dynamic regression relationships.
    \item Prewhitening.
    \item Separating intervention effects from ordinary temporal variation.
\end{itemize}

% ============================================================
\section{Module 11: Heteroscedasticity and Volatility Models}

\subsection{Chapter 12 -- Time Series Models of Heteroscedasticity}

\textbf{Major Topics}

\begin{itemize}
    \item Features of financial time series.
    \item ARCH models.
    \item GARCH models.
    \item Maximum likelihood estimation.
    \item Model diagnostics.
    \item Nonnegativity of conditional variances.
    \item Extensions of GARCH.
    \item Financial exchange-rate examples.
\end{itemize}

\subsection{ARCH Model}

A basic ARCH model may be expressed as

\[
X_t=\sigma_t\epsilon_t,
\]

with

\[
\sigma_t^2
=
\omega+
\alpha_1X_{t-1}^2.
\]

\subsection{GARCH Model}

A GARCH$(p,q)$ variance equation has the general form

\[
\sigma_t^2
=
\omega+
\sum_{i=1}^{q}\alpha_i X_{t-i}^2
+
\sum_{j=1}^{p}\beta_j\sigma_{t-j}^2.
\]

\textbf{Learning Objectives}

Students will be able to:

\begin{enumerate}
    \item Explain conditional heteroscedasticity.
    \item Recognize volatility clustering in financial time series.
    \item Construct an ARCH model.
    \item Construct and interpret GARCH models.
    \item Estimate volatility-model parameters.
    \item Diagnose fitted volatility models.
    \item Explain conditions required for nonnegative conditional variances.
    \item Understand extensions of the basic GARCH framework.
\end{enumerate}

\textbf{Difficult Topics}

\begin{itemize}
    \item Conditional variance dynamics.
    \item Maximum likelihood estimation for volatility models.
    \item Parameter restrictions ensuring nonnegative variances.
    \item Interpretation of ARCH/GARCH persistence.
\end{itemize}

% ============================================================
\section{Module 12: Spectral Analysis}

\subsection{Chapter 13 -- Introduction to Spectral Analysis}

\textbf{Major Topics}

\begin{itemize}
    \item Introduction to frequency-domain analysis.
    \item Periodogram.
    \item Spectral representation.
    \item Spectral distribution.
    \item Spectral density.
    \item Spectral densities for ARMA processes.
    \item Sampling properties of sample spectral density.
\end{itemize}

\textbf{Core Concept}

Spectral analysis studies how the variability of a time series is distributed
across frequencies rather than only across time lags.

A spectral density can be related to the autocovariance sequence through a
Fourier representation of the form

\[
f(\lambda)
=
\frac{1}{2\pi}
\sum_{k=-\infty}^{\infty}
\gamma(k)e^{-ik\lambda}.
\]

\textbf{Learning Objectives}

Students will be able to:

\begin{enumerate}
    \item Explain the motivation for frequency-domain analysis.
    \item Interpret the periodogram.
    \item Explain spectral representation.
    \item Define and interpret spectral density.
    \item Relate spectral density to autocovariance.
    \item Derive or interpret spectral densities for ARMA processes.
    \item Understand sampling properties of spectral estimates.
\end{enumerate}

\textbf{Difficult Topics}

\begin{itemize}
    \item Fourier representations.
    \item Frequency versus time-domain interpretations.
    \item Spectral density.
    \item Relationship between ACF and spectrum.
\end{itemize}

% ============================================================
\section{Module 13: Spectrum Estimation}

\subsection{Chapter 14 -- Estimating the Spectrum}

\textbf{Major Topics}

\begin{itemize}
    \item Smoothing spectral densities.
    \item Bias and variance.
    \item Bandwidth.
    \item Confidence intervals for the spectrum.
    \item Leakage.
    \item Tapering.
    \item Autoregressive spectrum estimation.
    \item Simulated-data examples.
    \item Actual-data examples.
    \item Other spectral-estimation methods.
\end{itemize}

\textbf{Learning Objectives}

Students will be able to:

\begin{enumerate}
    \item Explain why raw spectral estimates may be unstable.
    \item Understand spectral smoothing.
    \item Explain the bias--variance tradeoff.
    \item Interpret bandwidth.
    \item Construct or interpret confidence intervals for spectra.
    \item Explain spectral leakage.
    \item Understand tapering.
    \item Compare alternative methods of spectrum estimation.
\end{enumerate}

\textbf{Difficult Topics}

\begin{itemize}
    \item Bias--variance tradeoff in spectral estimation.
    \item Bandwidth selection.
    \item Leakage.
    \item Tapering.
    \item Statistical inference for spectra.
\end{itemize}

% ============================================================
\section{Module 14: Nonlinear Time-Series Models}

\subsection{Chapter 15 -- Threshold Models}

\textbf{Major Topics}

\begin{itemize}
    \item Graphical exploration of nonlinearity.
    \item Tests for nonlinearity.
    \item Polynomial models and explosiveness.
    \item First-order threshold autoregressive models.
    \item General threshold models.
    \item Tests for threshold nonlinearity.
    \item Estimation of TAR models.
    \item Model diagnostics.
    \item Prediction.
\end{itemize}

\subsection{Threshold Autoregressive Structure}

A simple threshold autoregressive model can be represented schematically as

\[
X_t=
\begin{cases}
c_1+\phi_{11}X_{t-1}+\epsilon_t,
& X_{t-d}\leq r,\\[4pt]
c_2+\phi_{21}X_{t-1}+\epsilon_t,
& X_{t-d}>r,
\end{cases}
\]

where $r$ is a threshold and $d$ is the delay parameter.

\textbf{Learning Objectives}

Students will be able to:

\begin{enumerate}
    \item Recognize evidence of nonlinear temporal dependence.
    \item Explain why polynomial models can become explosive.
    \item Define a threshold autoregressive model.
    \item Interpret regimes and threshold parameters.
    \item Test for threshold nonlinearity.
    \item Estimate TAR models.
    \item Diagnose fitted nonlinear models.
    \item Produce predictions from threshold models.
\end{enumerate}

\textbf{Difficult Topics}

\begin{itemize}
    \item Detecting nonlinearity.
    \item Threshold identification.
    \item Regime-dependent dynamics.
    \item Nonlinear estimation.
    \item Prediction under regime switching.
\end{itemize}

% ============================================================
\section{Module 15: Statistical Computing with R}

The supplied table of contents includes an extensive computational component,
including chapter-specific R commands, an R introduction, TSA library
functions, and dataset information.

\subsection{Computational Skills}

Students should develop the ability to:

\begin{enumerate}
    \item Import time-series datasets.
    \item Visualize time series.
    \item Calculate descriptive statistics.
    \item Examine ACF and PACF plots.
    \item Fit AR, MA, ARMA, and ARIMA models.
    \item Examine residuals.
    \item Generate forecasts.
    \item Fit seasonal models.
    \item Analyze interventions and outliers.
    \item Fit ARCH/GARCH-type models.
    \item Perform spectral analysis.
    \item Fit threshold models.
    \item Reproduce textbook examples computationally.
\end{enumerate}

\textbf{Recommended Computational Workflow}

\[
\boxed{
\text{Import Data}
\rightarrow
\text{Plot}
\rightarrow
\text{Transform}
\rightarrow
\text{Identify}
\rightarrow
\text{Fit}
\rightarrow
\text{Diagnose}
\rightarrow
\text{Forecast}
}
\]

% ============================================================
\section{Integrated Model-Building Workflow}

The entire course can be viewed as one unified modeling pipeline.

\subsection{Stage 1 -- Understand the Data}

\begin{itemize}
    \item Plot the series.
    \item Identify trend.
    \item Identify seasonality.
    \item Identify changing variance.
    \item Search for unusual observations.
\end{itemize}

\subsection{Stage 2 -- Determine Stationarity}

Ask:

\begin{itemize}
    \item Is the mean stable?
    \item Is the variance stable?
    \item Is the dependence structure stable?
    \item Is differencing required?
    \item Is a transformation required?
\end{itemize}

\subsection{Stage 3 -- Specify the Model}

Use:

\begin{itemize}
    \item ACF.
    \item PACF.
    \item Extended autocorrelation information.
    \item Graphical diagnostics.
    \item Subject-matter considerations.
\end{itemize}

Possible models include:

\[
\MA(q),\qquad
\AR(p),\qquad
\ARMA(p,q),\qquad
\ARIMA(p,d,q).
\]

For seasonal data:

\[
\text{Seasonal ARIMA}.
\]

For conditional variance dynamics:

\[
\ARCH,\qquad \GARCH.
\]

For nonlinear dynamics:

\[
\TAR.
\]

\subsection{Stage 4 -- Estimate Parameters}

Candidate approaches include:

\begin{itemize}
    \item Method of moments.
    \item Least squares.
    \item Maximum likelihood.
    \item Unconditional least squares.
    \item Bootstrap methods.
\end{itemize}

\subsection{Stage 5 -- Diagnose the Model}

Check:

\begin{itemize}
    \item Residual autocorrelation.
    \item Residual behavior.
    \item Remaining systematic structure.
    \item Parameter redundancy.
    \item Overfitting.
\end{itemize}

A satisfactory model should leave residuals that contain little predictable
structure.

\subsection{Stage 6 -- Forecast}

Generate:

\begin{itemize}
    \item One-step forecasts.
    \item Multi-step forecasts.
    \item Prediction limits.
    \item Updated forecasts.
\end{itemize}

\subsection{Stage 7 -- Compare and Refine}

The model-building process is iterative:

\[
\boxed{
\text{Specify}
\rightarrow
\text{Estimate}
\rightarrow
\text{Diagnose}
\rightarrow
\text{Revise}
\rightarrow
\text{Forecast}
}
\]

% ============================================================
\section{Major Difficult Topics Across the Course}

The following topics are likely to require the greatest conceptual and
mathematical effort.

\begin{longtable}{p{0.28\textwidth}p{0.62\textwidth}}
\toprule
\textbf{Topic} & \textbf{Why It Is Difficult} \\
\midrule
Stationarity &
Requires students to rethink ordinary statistical assumptions in terms of
dependence across time. \\

AR stationarity &
Requires understanding characteristic roots and stationarity regions. \\

Invertibility &
Requires distinguishing the representational properties of MA and ARMA
models. \\

ARIMA modeling &
Combines autoregression, moving averages, and differencing. \\

ACF/PACF identification &
Requires pattern recognition and understanding finite-sample behavior. \\

Parameter estimation &
Likelihood and least-squares methods become more complicated under temporal
dependence. \\

Residual diagnostics &
Students must assess whether temporal structure remains after fitting. \\

MMSE forecasting &
Requires conditional expectation and forecast-error reasoning. \\

Prediction limits &
Requires understanding uncertainty that accumulates with forecast horizon. \\

Seasonal ARIMA &
Combines seasonal and nonseasonal dynamics. \\

Spurious correlation &
Requires recognizing that strong relationships can arise from common
time-dependent structure rather than substantive relationships. \\

Prewhitening &
Requires understanding how temporal dependence can distort regression and
cross-correlation analysis. \\

ARCH/GARCH &
Introduces dynamic conditional variance in addition to dynamic conditional
means. \\

Spectral analysis &
Requires moving from the time domain to the frequency domain. \\

Spectral estimation &
Introduces smoothing, bandwidth, leakage, tapering, and bias--variance
tradeoffs. \\

Threshold models &
Requires nonlinear and regime-dependent thinking rather than a single global
linear model. \\
\bottomrule
\end{longtable}

% ============================================================
\section{Learning Objectives by Cognitive Level}

\subsection{Foundational Knowledge}

Students should be able to:

\begin{itemize}
    \item Define a stochastic process.
    \item Define stationarity.
    \item Define autocorrelation.
    \item Define AR, MA, ARMA, and ARIMA models.
    \item Define ARCH and GARCH models.
    \item Define spectral density.
    \item Define threshold models.
\end{itemize}

\subsection{Conceptual Understanding}

Students should be able to:

\begin{itemize}
    \item Explain temporal dependence.
    \item Explain stationarity.
    \item Explain differencing.
    \item Explain model identification.
    \item Explain residual diagnostics.
    \item Explain forecasting uncertainty.
    \item Explain conditional heteroscedasticity.
    \item Explain time-domain versus frequency-domain analysis.
    \item Explain nonlinear threshold behavior.
\end{itemize}

\subsection{Application}

Students should be able to:

\begin{itemize}
    \item Fit ARMA and ARIMA models.
    \item Analyze seasonal series.
    \item Fit regression models with time-series structure.
    \item Fit ARCH/GARCH models.
    \item Estimate spectra.
    \item Fit threshold models.
\end{itemize}

\subsection{Analysis}

Students should be able to:

\begin{itemize}
    \item Diagnose stationarity.
    \item Interpret ACF/PACF behavior.
    \item Compare competing models.
    \item Diagnose residual dependence.
    \item Evaluate forecasting performance.
    \item Analyze volatility dynamics.
    \item Detect nonlinear behavior.
\end{itemize}

\subsection{Evaluation}

Students should be able to:

\begin{itemize}
    \item Determine whether a model is adequate.
    \item Evaluate whether a model is overfitted.
    \item Compare alternative specifications.
    \item Evaluate forecast uncertainty.
    \item Assess whether a linear model is sufficient.
\end{itemize}

\subsection{Creation}

Students should ultimately be able to:

\begin{itemize}
    \item Develop a complete time-series model from raw data.
    \item Justify model choices statistically.
    \item Estimate and diagnose the selected model.
    \item Produce forecasts.
    \item Communicate model assumptions, limitations, and conclusions.
\end{itemize}

% ============================================================
\section{Suggested Course Sequence}

A coherent semester sequence is:

\begin{longtable}{p{0.10\textwidth}p{0.27\textwidth}p{0.48\textwidth}}
\toprule
\textbf{Weeks} & \textbf{Module} & \textbf{Content} \\
\midrule
1--2 & Foundations &
Introduction; stochastic processes; moments; covariance; stationarity \\

3 & Trends &
Deterministic and stochastic trends; regression; residual analysis \\

4--5 & Stationary Models &
Linear processes; MA; AR; ARMA; invertibility \\

6 & Nonstationary Models &
Differencing; ARIMA; transformations \\

7 & Specification &
ACF; PACF; model identification; nonstationarity \\

8 & Estimation &
Moments; least squares; maximum likelihood; bootstrap \\

9 & Diagnostics &
Residual analysis; overfitting; parameter redundancy \\

10 & Forecasting &
MMSE forecasts; ARIMA forecasting; prediction limits; updating \\

11 & Seasonal Models &
Seasonal ARIMA; multiplicative models; seasonal forecasting \\

12 & Regression &
Interventions; outliers; spurious correlation; prewhitening \\

13 & Heteroscedasticity &
Financial series; ARCH; GARCH; diagnostics \\

14 & Spectral Analysis &
Periodogram; spectral density; spectrum estimation \\

15 & Nonlinear Models &
Nonlinearity; threshold models; TAR estimation; diagnostics; prediction \\

16 & Integration &
Comprehensive modeling project and computational review \\
\bottomrule
\end{longtable}

% ============================================================
\section{Recommended Assessment Structure}

A course based on this structure can use the following assessment framework.

\begin{longtable}{p{0.45\textwidth}p{0.20\textwidth}p{0.25\textwidth}}
\toprule
\textbf{Assessment} & \textbf{Weight} & \textbf{Primary Skills} \\
\midrule
Problem Sets & 15\% &
Theory, derivations, calculations \\

Computational Laboratories & 15\% &
R implementation and interpretation \\

Midterm Examination & 20\% &
Foundations through model specification \\

Modeling Project & 20\% &
Complete empirical modeling workflow \\

Final Examination & 25\% &
Theory, modeling, diagnostics, forecasting \\

Participation / Presentation & 5\% &
Communication and interpretation \\
\bottomrule
\end{longtable}

% ============================================================
\section{Suggested Capstone Project}

Students should complete a full empirical time-series analysis following the
course model-building strategy.

\subsection{Project Requirements}

\begin{enumerate}
    \item Select an appropriate real-world time series.
    \item Describe the data-generating context.
    \item Produce an exploratory time-series plot.
    \item Assess trend and seasonality.
    \item Investigate stationarity.
    \item Apply transformations or differencing when appropriate.
    \item Examine ACF and PACF.
    \item Specify candidate models.
    \item Estimate model parameters.
    \item Perform residual diagnostics.
    \item Compare competing models.
    \item Select a final model.
    \item Generate forecasts.
    \item Quantify forecast uncertainty.
    \item Interpret the results in substantive terms.
    \item Document the analysis using reproducible R code.
\end{enumerate}

\subsection{Advanced Project Extensions}

Students seeking additional difficulty may extend the project to include:

\begin{itemize}
    \item Seasonal ARIMA.
    \item Intervention analysis.
    \item Time-series regression.
    \item ARCH/GARCH modeling.
    \item Spectral analysis.
    \item Threshold autoregressive modeling.
\end{itemize}

% ============================================================
\section{Concept Dependency Structure}

The conceptual dependencies of the course can be summarized as follows:

\[
\begin{array}{c}
\text{Probability and Statistics}\\
\downarrow\\
\text{Stochastic Processes}\\
\downarrow\\
\text{Stationarity}\\
\downarrow\\
\text{ACF/PACF and Dependence}\\
\downarrow\\
\text{AR / MA Models}\\
\downarrow\\
\text{ARMA Models}\\
\downarrow\\
\text{Differencing}\\
\downarrow\\
\text{ARIMA Models}\\
\downarrow\\
\text{Model Specification}\\
\downarrow\\
\text{Parameter Estimation}\\
\downarrow\\
\text{Model Diagnostics}\\
\downarrow\\
\text{Forecasting}
\end{array}
\]

After the core ARIMA framework is established, the course branches into:

\[
\begin{array}{ccc}
\text{ARIMA Core}
&\longrightarrow&
\text{Seasonal Models}\\
&\longrightarrow&
\text{Time-Series Regression}\\
&\longrightarrow&
\text{ARCH/GARCH}\\
&\longrightarrow&
\text{Spectral Analysis}\\
&\longrightarrow&
\text{Threshold Models}.
\end{array}
\]

% ============================================================
\section{Final Course Competency Map}

At the completion of the course, a successful student should be capable of
performing the following complete sequence:

\[
\boxed{
\begin{minipage}{0.88\textwidth}
\centering
\textbf{Understand the data}
$\rightarrow$
\textbf{Identify temporal structure}
$\rightarrow$
\textbf{Assess stationarity}
$\rightarrow$
\textbf{Specify a model}
$\rightarrow$
\textbf{Estimate parameters}
$\rightarrow$
\textbf{Diagnose the model}
$\rightarrow$
\textbf{Revise if necessary}
$\rightarrow$
\textbf{Forecast}
$\rightarrow$
\textbf{Quantify uncertainty}
$\rightarrow$
\textbf{Communicate conclusions}
\end{minipage}
}
\]

\subsection{Core Competencies}

\begin{enumerate}
    \item \textbf{Probability and stochastic-process reasoning}
    \item \textbf{Stationarity and temporal dependence}
    \item \textbf{AR/MA/ARMA modeling}
    \item \textbf{ARIMA and seasonal ARIMA modeling}
    \item \textbf{Model identification using ACF/PACF}
    \item \textbf{Parameter estimation}
    \item \textbf{Residual diagnostics}
    \item \textbf{Forecasting and prediction intervals}
    \item \textbf{Dynamic regression}
    \item \textbf{Volatility modeling}
    \item \textbf{Frequency-domain analysis}
    \item \textbf{Nonlinear threshold modeling}
    \item \textbf{Statistical programming and reproducibility}
\end{enumerate}

% ============================================================
\section{Textbook-to-Course Mapping}

The supplied contents explicitly organize the material into Chapters 1--15,
followed by R commands, datasets, bibliography, and index. The first four
chapters establish introduction, fundamental concepts, trends, and stationary
models; Chapters 5--9 develop nonstationary models, specification, estimation,
diagnostics, and forecasting. 

The later chapters extend the framework to seasonal models, time-series
regression, heteroscedasticity, spectral analysis, spectrum estimation, and
threshold models.

\begin{longtable}{p{0.12\textwidth}p{0.35\textwidth}p{0.38\textwidth}}
\toprule
\textbf{Chapter} & \textbf{Textbook Topic} & \textbf{Course Role} \\
\midrule
1 & Introduction & Orientation and model-building strategy \\
2 & Fundamental Concepts & Mathematical foundation \\
3 & Trends & Trend modeling and regression \\
4 & Stationary Time Series & AR, MA, ARMA foundations \\
5 & Nonstationary Time Series & Differencing and ARIMA \\
6 & Model Specification & Identification \\
7 & Parameter Estimation & Statistical estimation \\
8 & Model Diagnostics & Model validation \\
9 & Forecasting & Prediction and uncertainty \\
10 & Seasonal Models & Seasonal extensions \\
11 & Time Series Regression & Interventions and dynamic regression \\
12 & Heteroscedasticity & ARCH/GARCH volatility modeling \\
13 & Spectral Analysis & Frequency-domain foundations \\
14 & Estimating the Spectrum & Statistical spectral estimation \\
15 & Threshold Models & Nonlinear time-series analysis \\
\bottomrule
\end{longtable}

% ============================================================
\section{Conclusion}

This course structure turns the supplied textbook contents into a coherent
progression from fundamental probability concepts to advanced time-series
modeling.

The central intellectual progression is:

\[
\boxed{
\text{Describe}
\rightarrow
\text{Understand}
\rightarrow
\text{Model}
\rightarrow
\text{Estimate}
\rightarrow
\text{Diagnose}
\rightarrow
\text{Forecast}
}
\]

The first half of the course establishes the linear time-domain framework,
with AR, MA, ARMA, ARIMA, identification, estimation, diagnostics, and
forecasting as its central components. The second half broadens the framework
to seasonal dependence, dynamic regression, conditional heteroscedasticity,
frequency-domain methods, and nonlinear threshold dynamics.

This organization therefore supports both a mathematically rigorous course
and a practical empirical modeling course in which students repeatedly apply
the same model-building logic to increasingly sophisticated classes of
time-series models.

\end{document}
```
